\chapter{Pendahuluan}

\section{Latar Belakang}

Hadis adalah segala sesuatu yang disandarkan kepada Nabi Muhammad saw., baik berupa ucapan, perbuatan, persetujuan (\textit{taqrir}), maupun hal-hal lain yang berkaitan dengan beliau \cite{handayani2026hadis}.
Sebagai sumber hukum Islam kedua setelah Al-Qur'an, hadis berperan penting dalam menjelaskan, merinci, dan memperkuat ketentuan-ketentuan syariat \cite{hatta2026teknik}.
Secara struktural, hadis terdiri atas dua komponen utama: \textit{matan}, yaitu isi atau teks hadis itu sendiri, dan \textit{sanad}, yaitu rangkaian perawi yang menyampaikan matan tersebut secara berkesinambungan hingga sampai kepada Nabi Muhammad saw.
Dengan kedudukan hadis sebagai pedoman utama umat Islam, menjaga keaslian dan kesahihan kedua komponen ini menjadi suatu keharusan.

Sanad, sebagai rantai transmisi hadis, pada hakikatnya merupakan jaringan keilmuan Islam yang terbentuk dari hubungan guru–murid yang berlangsung secara berkesinambungan lintas generasi \cite{hatta2026teknik}.
Setiap individu yang terlibat dalam rantai ini disebut perawi; mereka adalah ulama dan cendekiawan Muslim yang menjadi mata rantai dalam pewarisan ilmu keagamaan dari satu generasi ke generasi berikutnya.
Identifikasi seorang perawi dalam jaringan ini umumnya dilakukan dengan menelusuri daftar guru yang meriwayatkan hadis kepadanya serta daftar murid yang mengambil hadis darinya, di samping memperhatikan generasi (\textit{tabaqat}) dan masa hidupnya \cite{WanKamalNadzif2025135}.
Karena strukturnya yang relasional inilah, sanad sangat tepat direpresentasikan sebagai jaringan berarah, dengan perawi sebagai simpul (\textit{node}) dan hubungan guru–murid sebagai sisi (\textit{edge}) yang menghubungkan antarsimpul \cite{FAROOQI2024110439}.

Meskipun seluruh perawi terhubung dalam jaringan periwayatan, tidak semua perawi memiliki peran yang setara. Sebagian perawi menempati posisi strategis karena menerima hadis dari banyak guru sekaligus menyebarkannya kepada banyak murid, sehingga menjadi titik temu yang dilalui beragam jalur sanad. Perawi semacam ini dikenal sebagai \textit{common link}, dan keberadaannya menegaskan pentingnya mengidentifikasi perawi yang paling berpengaruh serta paling banyak menghubungkan antarjalur periwayatan \cite{mufid2023track}.
Pada kajian klasik, identifikasi tersebut umumnya dilakukan secara manual melalui penelusuran \textit{isnad bundle}, tetapi ketika cakupan kajian melibatkan ratusan ribu perawi dengan relasi yang sangat kompleks, pendekatan manual menjadi tidak memadai.
Kondisi inilah yang mendorong kebutuhan akan analisis kuantitatif dan komputasional untuk mengukur posisi serta pengaruh tiap perawi secara sistematis dan terukur dalam skala besar \cite{FAROOQI2024110439}.

Kebutuhan akan pendekatan komputasional tersebut sesungguhnya telah mulai dijawab sejak lebih dari satu dekade lalu melalui sejumlah penelitian yang menganalisis jaringan perawi hadis secara kuantitatif. Beberapa penelitian membangun jaringan sosial historis perawi lintas generasi dan menerapkan analisis jaringan untuk merekonstruksi informasi temporal dari struktur transmisinya \cite{Shuaib2022}, 
sementara penelitian lain merepresentasikan jaringan perawi dalam bentuk \textit{knowledge graph} dan menganalisisnya menggunakan berbagai perangkat lunak jaringan \cite{FAROOQI2024110439}. 
Sebagaimana ditinjau oleh Mghari \textit{et al}, sejumlah studi bahkan telah menerapkan metrik sentralitas seperti \textit{in-degree}, \textit{out-degree}, \textit{betweenness}, dan \textit{closeness} untuk mengidentifikasi perawi paling berpengaruh dalam jaringan periwayatan, dan upaya ke arah visualisasi yang lebih eksploratif pun mulai dirintis melalui prototipe seperti Chain of Hadith Narrators Visualizer (CHN) \cite{Mghari2024Narrator2Vec} .
Penelitian-penelitian ini menegaskan bahwa jaringan perawi memang dapat dimodelkan dan dianalisis secara komputasional dalam skala besar. Meskipun demikian, luaran yang dihasilkan umumnya bersifat statis dan bergantung pada perangkat lunak desktop khusus seperti Gephi dan Cytoscape, sehingga sulit diakses oleh kalangan yang lebih luas \cite{FAROOQI2024110439}. Hasil analisis yang diperoleh pun belum disajikan dalam bentuk yang dapat dieksplorasi secara interaktif oleh pelajar, peneliti, maupun masyarakat umum.

Meski demikian, representasi \textit{knowledge graph} yang telah digunakan dalam penelitian - penelitian tersebut sesungguhnya menawarkan fondasi yang dapat dikembangkan menjadi sistem yang lebih interaktif dan mudah diakses secara luas.
\textit{Knowledge graph} menggambarkan entitas beserta relasi antarentitas dalam bentuk graf, sehingga keterhubungan data yang kompleks dapat disajikan secara semantis dan terstruktur, bukan sekadar sebagai data tabular yang terpisah.
Karena keunggulan tersebut, \textit{knowledge graph} telah diadopsi secara luas di berbagai bidang, seperti pendidikan untuk memprediksi kinerja siswa maupun pariwisata untuk menggali wawasan dari data berskala besar \cite{Aoula2026}.
Dalam kajian hadis, pendekatan ini menyediakan representasi yang sesuai untuk memodelkan jaringan periwayatan secara komputasional, dengan perawi sebagai entitas dan hubungan guru–murid sebagai relasi yang menghubungkan antarsimpul \cite{FAROOQI2024110439}.

Setelah representasi berbasis \textit{knowledge graph} ditetapkan, komponen yang diperlukan berikutnya adalah sumber data sebagai bahan pembentukan jaringan perawi. Penelitian ini menggunakan Sanadset 650K, sebuah dataset sanad hadis yang dipublikasikan secara terbuka oleh Mghari \textit{et al}.
Melalui aktivitas kurasi yang komprehensif, mereka menghimpun 650.986 entri hadis yang diekstraksi dari 926 kitab hadis klasik berbahasa Arab, dengan setiap entri menandai komponen sanad, tiap perawi di dalamnya, serta matan secara terstruktur. Pemilihan dataset ini didasarkan pada keautentikannya karena bersumber langsung dari kitab-kitab hadis klasik, cakupannya yang luas lintas ratusan kitab, serta skalanya yang besar sehingga mampu merepresentasikan jaringan periwayatan secara komprehensif \cite{MGHARI2022108540}.
Kredibilitas Sanadset 650K turut ditunjukkan oleh pemanfaatannya dalam penelitian lanjutan tentang representasi perawi \cite{Mghari2024Narrator2Vec}. Dengan data yang autentik dan berskala besar inilah jaringan perawi yang akurat dan berbasis bukti dapat dibangun dalam penelitian ini.

Dengan jaringan perawi yang telah terbentuk, dibutuhkan metode untuk mengukur posisi dan pengaruh tiap perawi di dalamnya.
\textit{Social Network Analysis} (SNA) merupakan pendekatan yang sesuai untuk tujuan tersebut, dengan memanfaatkan sejumlah metrik sentralitas untuk mengukur peran tiap simpul (\textit{node}) dalam jaringan, di antaranya \textit{degree centrality (in-degree dan out-degree)}, \textit{betweenness}, \textit{closeness}, dan \textit{eigenvector centrality} \cite{LópezWerner04072025}.
Melalui metrik-metrik tersebut, perawi yang paling berpengaruh maupun yang paling berperan sebagai penghubung antarjalur periwayatan dapat diidentifikasi secara sistematis dan terukur—sebagaimana telah dibuktikan dalam sejumlah kajian kuantitatif jaringan perawi yang telah diulas sebelumnya \cite{FAROOQI2024110439}.

Jaringan perawi beserta hasil analisis sentralitasnya hanya akan bermakna apabila dapat ditelusuri dengan mudah, sehingga diperlukan antarmuka (\textit{front-end}) interaktif berbasis web sebagai sarana penyajiannya. Kebutuhan ini didasarkan pada tiga pertimbangan. Pertama, dari sisi interaktivitas, keterhubungan antarperawi tidak cukup disajikan sebagai gambar statis; pengguna memerlukan kemampuan untuk memperbesar, menyaring, mencari, dan menelusuri tiap simpul secara langsung.
Kedua, dari sisi aksesibilitas, antarmuka berbasis web dapat diakses tanpa instalasi perangkat lunak khusus dan menjangkau khalayak luas, mulai dari pelajar, peneliti, hingga masyarakat umum. Ketiga, dari sisi penyajian, menampilkan keseluruhan jaringan yang sangat besar sekaligus justru mengurangi keterbacaan; sebagaimana ditempuh dalam analisis jaringan lainnya, hasil pengukuran sentralitas dapat dimanfaatkan untuk menonjolkan simpul-simpul paling sentral sebagai fokus visualisasi \cite{LópezWerner04072025}.
Ketiga pertimbangan inilah yang mengarahkan penelitian ini pada pengembangan \textit{front-end} interaktif berbasis web untuk \textit{knowledge graph} perawi.

\section{Rumusan Masalah}

Selama ini penelitian terkait sanad hadis umumnya hanya berhenti pada tahap visualisasi hubungan antarperawi tanpa mengkaji lebih dalam peran dan pengaruh masing-masing perawi dalam jaringan tersebut. Selain itu, belum banyak penelitian yang menyajikan hasil analisis jaringan tersebut dalam bentuk \textit{front-end} interaktif yang memudahkan eksplorasi jaringan keilmuan islam secara langsung. Berdasarkan hal tersebut, rumusan masalah dalam penelitian ini adalah sebagai berikut:


\begin{enumerate}
	\item Bagaimana membangun \textit{Knowledge Graph} jaringan keilmuan Islam berdasarkan relasi guru-murid yang diperoleh dari data sanad hadis?
	\item Bagaimana mengidentifikasi peran dan pengaruh masing-masing perawi dalam jaringan sanad hadis?
	\item Bagaimana membangun sistem visualisasi \textit{front-end} yang mampu menyajikan jaringan perawi dan hasil analisis jaringan keilmuan Islam secara interaktif?
\end{enumerate}

\section{Tujuan Penelitian}

Penelitian ini bertujuan untuk mengembangkan Front-End Knowledge Graph Keagamaan Ilmuwan Islam dengan memanfaatkan data sanad hadis dan pendekatan \textit{Social Network Analysis} (SNA). Secara khusus, tujuan penelitian ini adalah:

\begin{enumerate}
	\item Menghasilkan representasi \textit{Knowledge Graph} yang menggambarkan hubungan guru dan murid antarperawi berdasarkan data sanad hadis.
	\item Menganalisis jaringan keilmuan Islam menggunakan metode \textit{Social Network Analysis} (SNA) melalui perhitungan metrik sentralitas untuk mengidentifikasi perawi yang memiliki peran penting dalam jaringan.
	\item Membangun sistem visualisasi \textit{front-end} yang mampu menyajikan jaringan perawi dan hasil analisis jaringan keilmuan Islam secara interaktif dan mudah digunakan.
\end{enumerate}	

\section{Batasan Penelitian}

Batasan penelitian dalam skripsi ini adalah sebagai berikut:

\begin{enumerate}
	\item Penelitian ini menggunakan dataset publik \textit{Sanadset 650K: Data on Hadith Narrators} sebagai sumber data utama. Pengolahan data hanya dilakukan pada kolom sanad yang tersedia dalam dataset tersebut, tidak mencakup kolom lainnya seperti matan, buku, dan Hadis.

	\item Nama perawi di kolom sanad digunakan sebagaimana tercantum dalam dataset tanpa normalisasi atau verifikasi keaslian nama aslinya. Jaringan yang dibangun merepresentasikan hubungan antar perawi tanpa memperhitungkan kronologi periwayatan maupun isi matan hadis.

	\item Proses pembangunan \textit{Knowledge Graph} didasarkan pada relasi guru-murid yang diekstraksi dari data sanad, sehingga entitas yang dimodelkan hanya mencakup perawi hadis beserta hubungan periwayatannya.

	\item Analisis jaringan menggunakan metode \textit{Social Network Analysis} (SNA) dengan menggunakan 5 metrik sentralitas, yaitu \textit{in degree centrality}, \textit{out degree centrality}, \textit{eigenvector centrality}, \textit{closeness centrality} dan \textit{betweenness centrality}.

	\item Pengembangan sistem dibatasi pada antarmuka berbasis web (\textit{web-based front-end}) yang mencakup visualisasi graf interaktif dan eksplorasi jaringan keilmuan Islam.

	\item Penelitian dilaksanakan pada tahun 2026 dengan menggunakan dataset publik \textit{Sanadset 650K} serta perangkat keras pribadi penulis sebagai sarana komputasi.
\end{enumerate}


%\noindent \textbf{Contoh penulisan batasan skripsi Teknik Elektro:}
%
%%-------------------------------------------------	
%\noindent\fbox{%
%	\parbox{\textwidth}{%
%\begin{enumerate}
%\item Objek penelitian: Studi efisiensi dan kinerja sistem pemantauan suhu dan arus pada sistem pembangkit tenaga listrik.
%\item Metode penelitian: Penelitian eksperimental menggunakan analisis simulasi dan pengujian sistem pada skala laboratorium.
%\item Waktu dan tempat penelitian: Waktu penelitian adalah Maret-Agustus 2022 di 
%FasLab TTL.
%\item Populasi dan sampel: Populasi adalah sistem pembangkit tenaga listrik, dan sampel diambil sebanyak 3 sistem yang berbeda.
%\item Variabel: Variabel bebas adalah metode pemantauan suhu dan arus, dan variabel terikat adalah efisiensi dan kinerja
%\item Hipotesis: bahwa metode pemantauan suhu dan arus berpengaruh terhadap efisiensi dan kinerja sistem pembangkit tenaga listrik.
%\item Keterbatasan Penelitian: Keterbatasan penelitian adalah hanya melibatkan 
%pengujian sistem pada skala laboratorium dan membatasi analisis pada variabel 
%bebas dan terikat.
%\end{enumerate}
%		
%	}%
%}
%
%%-------------------------------------------------	
%\newpage
%
%\noindent \textbf{Contoh penulisan batasan skripsi Teknik Biomedis:}
%
%%-------------------------------------------------	
%\noindent\fbox{%
%	\parbox{\textwidth}{%
%\begin{enumerate}
%\item Objek Penelitian: Studi efektivitas dan keamanan alat elektromedik seperti pacu jantung.
%\item Metode Penelitian: Penelitian deskriptif dan observasional menggunakan survei dan analisis data.
%\item Waktu dan Tempat Penelitian: Waktu penelitian adalah Januari-Juni 2022 di 
%Rumah Sakit Sarjito. 
%\item Populasi dan Sampel: Populasi adalah pasien yang menggunakan pacu jantung, dan sampel diambil dengan metode random sampling sebanyak 100 pasien. 
%\item Variabel: Variabel bebas nya adalah jenis alat pacu jantung, dan variabel terikatnya adalah efektivitas dan keamanan. 
%\item Hipotesis: bahwa jenis alat pacu jantung berpengaruh terhadap efektivitas dan keamanan. 
%\item Keterbatasan Penelitian: Keterbatasan penelitian adalah hanya melibatkan satu rumah sakit sebagai lokasi penelitian dan membatasi dalam analisis data hanya pada variabel bebas dan terikat.
%\end{enumerate}
%		
%	}%
%}
%
%%-------------------------------------------------	
%
%\newpage
%\noindent \textbf{Contoh penulisan batasan skripsi Teknologi Informasi:}
%
%%-------------------------------------------------	
%\noindent\fbox{%
%	\parbox{\textwidth}{%
%\begin{enumerate}
%\item Objek Penelitian: Analisis perbandingan efektivitas dan efisiensi antara sistem manajemen proyek tradisional dan sistem manajemen proyek berbasis teknologi informasi. 
%\item Metode Penelitian: Penelitian kualitatif dengan menggunakan wawancara dan 
%survei terhadap para pelaku proyek di berbagai perusahaan. 
%\item Waktu dan Tempat Penelitian: Waktu penelitian adalah Januari-Juni 2022 di 
%perusahaan-perusahaan di wilayah Bantul. 
%\item Populasi dan Sampel: Populasi nya adalah perusahaan yang melakukan proyek, dan sampel diambil sebanyak 10 perusahaan yang menerapkan sistem manajemen 
%proyek tradisional dan 10 perusahaan yang menggunakan sistem manajemen 
%proyek berbasis teknologi informasi. 
%\item Variabel: Variabel bebasnya adalah sistem manajemen proyek, dan variabel 
%terikatnya adalah efektivitas dan efisiensi. 
%\item Hipotesis: bahwa sistem manajemen proyek berbasis teknologi informasi lebih efektif dan efisien dibandingkan dengan sistem manajemen proyek tradisional.
%\item Keterbatasan Penelitian: Keterbatasan penelitian adalah penelitian hanya dilakukan pada perusahaan di wilayah Bantul dan hanya melibatkan wawancara dan survei sebagai metode pengumpulan data.
%\end{enumerate}
%		
%	}%
%}

\section{Manfaat Penelitian}

Penelitian ini diharapkan dapat memberikan manfaat bagi berbagai pihak sebagai berikut:

\begin{enumerate}
	\item \textbf{Bagi Peneliti}\\
	Penelitian ini meningkatkan pemahaman dan keterampilan peneliti dalam menerapkan pendekatan komputasional pada data keagamaan, khususnya dalam pengolahan data sanad hadis, pembangunan \textit{Knowledge Graph}, serta analisis jaringan menggunakan metode \textit{Social Network Analysis} (SNA). Pengalaman ini menjadi bekal bagi pengembangan penelitian atau proyek ilmiah berikutnya di bidang yang serupa.

	\item \textbf{Bagi Akademik}\\
	Penelitian ini memberikan kontribusi pada pengembangan ilmu di bidang teknologi informasi yang diterapkan pada domain keislaman. Hasil penelitian dapat menjadi referensi metodologis bagi penelitian selanjutnya yang ingin mengkaji jaringan keilmuan Islam melalui pendekatan \textit{Knowledge Graph} dan analisis jaringan secara kuantitatif. Selain itu, data perawi paling sentral yang dihasilkan dapat dibandingkan dengan penilaian ulama dalam literatur ilmu hadis, sehingga membuka peluang untuk menemukan perspektif baru dalam kajian hadis.

	\item \textbf{Bagi Mahasiswa dan Praktisi Ilmu Hadis}\\
	Sistem visualisasi \textit{front-end} yang dihasilkan memudahkan mahasiswa maupun praktisi ilmu hadis dalam menelusuri posisi dan otoritas seorang perawi dalam jaringan sanad. Pengguna dapat melihat siapa saja guru dan murid yang berinteraksi dengan seorang perawi, serta memahami pola transmisi keilmuan antar generasi. Hal ini memperkaya pemahaman terhadap sanad tidak hanya sebagai rantai teks, tetapi juga sebagai jaringan sosial keilmuan yang kompleks.

	\item \textbf{Bagi Masyarakat Umum}\\
	Masyarakat umum yang tertarik mempelajari sejarah keilmuan Islam dapat memanfaatkan sistem ini untuk mengeksplorasi jaringan perawi hadis secara mandiri dan interaktif, tanpa memerlukan pemahaman teknis yang mendalam. Dengan tampilan visual yang intuitif, informasi mengenai jaringan keilmuan Islam menjadi lebih mudah diakses dan dipahami oleh berbagai kalangan.
\end{enumerate}



\section{Sistematika Penulisan}

\noindent Bab I membahas pendahuluan yang mencakup latar belakang, rumusan masalah, tujuan penelitian, batasan penelitian, manfaat penelitian, dan sistematika penulisan.

\noindent Bab II membahas tinjauan pustaka dan dasar teori yang menjadi landasan penelitian. Tinjauan pustaka berisi kajian terhadap penelitian-penelitian terdahulu yang relevan dengan topik \textit{Knowledge Graph}, analisis jaringan sanad hadis, dan visualisasi data keagamaan. Dasar teori mencakup konsep hadis dan sanad, \textit{Knowledge Graph}, \textit{Social Network Analysis} beserta metrik sentralitasnya, serta teknologi pengembangan \textit{front-end} yang digunakan.

\noindent Bab III membahas metode penelitian yang digunakan, meliputi alat dan bahan penelitian, tahapan pengolahan data sanad dari dataset \textit{Sanadset 650K}, proses pembangunan \textit{Knowledge Graph} berbasis relasi guru-murid antarperawi, penerapan metode \textit{Social Network Analysis} untuk menganalisis jaringan perawi, serta perancangan dan pengembangan sistem visualisasi \textit{front-end} yang interaktif.

\noindent Bab IV membahas hasil dan pembahasan dari setiap tahapan penelitian. Bab ini menyajikan hasil pembangunan \textit{Knowledge Graph} jaringan perawi, hasil analisis SNA berupa nilai metrik sentralitas masing-masing perawi beserta interpretasinya, serta hasil pengembangan sistem visualisasi \textit{front-end} beserta pembahasannya.

\noindent Bab V membahas kesimpulan yang menjawab seluruh rumusan masalah penelitian berdasarkan hasil yang telah diperoleh, serta saran untuk pengembangan penelitian selanjutnya di bidang yang relevan.

